\section*{Reading notes}
\subsection*{\texttt{GELFAND-P001485}}
\par\noindent\leavevmode\phantomsection\label{GELFAND-0161-U001}\label{GELFAND-P001485:eq10-cancellation}\textbf{\texttt{GELFAND-0161-U001}} (eq10-cancellation).\nopagebreak[4]\par
In the exponent in (10), after the second integral and above the left edge of the surviving m, a small sign is densely cancelled. The surviving entry reads m lambda dt; below it, lambda squared = x-dot squared / m squared is explicitly written.\par
Candidates: a cancelled m; another cancelled single sign. Unresolved in the diplomatic source.\par
Historical source review: Original coordinates (1805,60,1870,126); wider formula (70,35,2210,230); the following stationary relation (1950,310,2230,475). Frame 2408 by 3480. The complete original, context at 3x and cancelled sign at 7x were examined. The strokes conceal the cancelled symbol. Formal differentiation of the surviving lambda\textasciicircum{}(-1) x-dot\textasciicircum{}2 + m lambda would give denominator m in the stationary relation, whereas the source writes m\textasciicircum{}2. Both source expressions were retained; neither m\textasciicircum{}2 nor a numerical coefficient was inserted into (10).\par
Current source comparison: The cancellation is still opaque at 5x/10x. The fresh comparison confirms surviving m and the separate m\textasciicircum{}2 denominator. The tool rescaled one wide 3x context; two narrower overlapping 3x views were then inspected and their exact native bytes verified.\par\medskip
\par\noindent\leavevmode\phantomsection\label{GELFAND-0161-U002}\label{GELFAND-P001485:eq11-differential}\textbf{\texttt{GELFAND-0161-U002}} (eq11-differential).\nopagebreak[4]\par
The last differential in the exponent in (11) has a distinctly x-shaped last letter; the t in neighboring formula (10) is written differently.\par
Candidates: dx; dt. Unresolved in the diplomatic source.\par
Historical source review: Original coordinates (980,460,1090,540). Frame 2408 by 3480. The complete original, formula at 3x and differential at 7x were examined. Visually dx is the stronger reading and was reproduced. The familiar expression using a time parameter would suggest dt, but that was not substituted for the visible entry.\par
Current source comparison: The x shape still favors dx at 5x/10x. The alternative dt is retained in the apparatus without repairing the formula.\par\medskip
\par\noindent\leavevmode\phantomsection\label{GELFAND-0161-U003}\label{GELFAND-P001485:metric-abbreviation}\textbf{\texttt{GELFAND-0161-U003}} (metric-abbreviation).\nopagebreak[4]\par
In parentheses after “one-dim. metric” the text reads “everything else ... from reparametrization”. The abbreviation before “from” resembles “свод.” or “своб.”; its last letter has a loop.\par
Candidates: свод.; своб.. Unresolved in the diplomatic source.\par
Historical source review: Original coordinates (1970,764,2120,860); wider parenthesis (1550,755,2270,930). Frame 2408 by 3480. Context at 3x, word at 7x and the word with the following “из” at 5x were examined. The last handwritten letter resembles Cyrillic д, but cannot confidently be distinguished from б in this hand; the abbreviation was not expanded.\par
Current source comparison: The loop still favors д without excluding б. The English body keeps the source abbreviation and does not supply an expansion.\par\medskip
\par\noindent\leavevmode\phantomsection\label{GELFAND-0161-U004}\label{GELFAND-P001485:eq12-identity}\textbf{\texttt{GELFAND-0161-U004}} (eq12-identity).\nopagebreak[4]\par
Formula (12) is clearly written as integral (d lambda / lambda) exp(a lambda) exp(b/lambda) = exp(sqrt(ab)). No integration limits, additional factors or coefficient 2 before the square root are visible.\par
Candidates: \textbackslash{}int\textbackslash{}frac\{d\textbackslash{}lambda\}\{\textbackslash{}lambda\}e\textasciicircum{}\{a\textbackslash{}lambda\}e\textasciicircum{}\{b/\textbackslash{}lambda\}=e\textasciicircum{}\{\textbackslash{}sqrt\{ab\}\}. Unresolved in the diplomatic source.\par
Historical source review: Original coordinates (290,864,1200,1070). Frame 2408 by 3480. The original and fragment at 3x were examined. The equality was retained as written. Its analytic meaning, convergence conditions and normalization are not established by the source; no mathematically expected factor or different function was inserted.\par
Current source comparison: The initial detail boxes clipped the bottom of the measure and the top of the radical. Corrected 5x/10x details show the whole denominator lambda and radical bar. They confirm the literal expression; a conditional convergent real integral is calculated separately.\par\medskip
\par\noindent\leavevmode\phantomsection\label{GELFAND-0161-U005}\label{GELFAND-P001485:eq14-differentials}\textbf{\texttt{GELFAND-0161-U005}} (eq14-differentials).\nopagebreak[4]\par
In (14), after the high-written integral G\_mu\_nu(x) sqrt(g) g\textasciicircum{}ab partial\_a x\textasciicircum{}mu partial\_b x\textasciicircum{}nu, d sigma dt is placed to the right on a lower line, approximately level with the main exponential e.\par
Candidates: d sigma dt belongs to the integral in the exponent; d sigma dt is written outside the exponent. Unresolved in the diplomatic source.\par
Historical source review: Original coordinates (850,1790,1840,2020); differentials (1640,1920,1840,2020). Frame 2408 by 3480. The original and fragment at 3x were examined. The transcription preserves the visible low placement on a second line of an aligned formula. Mathematical context more likely assigns the differentials to the integral, but the source layout alone does not finally establish the scope of the exponent.\par
Current source comparison: The fresh 5x/10x details confirm the low placement. They do not settle mathematical scope; the aligned transcription remains literal.\par\medskip
\subsection*{\texttt{GELFAND-P001486}}
\par\noindent\leavevmode\phantomsection\label{GELFAND-0161-U006}\label{GELFAND-P001486:feynman-verb}\textbf{\texttt{GELFAND-0161-U006}} (feynman-verb).\nopagebreak[4]\par
In I.M.G.’s remark, between “and” and “aesthetics”, a short joined entry “учи...” is visible; its ending is cramped.\par
Candidates: учить [to teach]; учите [teach; imperative plural or formal]; учиться [to learn]. Unresolved in the diplomatic source.\par
Historical source review: Original coordinates (950,2110,1145,2200). Frame 2408 by 3480. The complete original, line at 3x and fragment at 7x were examined. Visually “учить” is preferred; “учиться” is only a possible reading of the cramped ending strokes and was not accepted on grammatical grounds.\par
Current source comparison: The cramped ending remains unresolved at 5x/10x. The preferred historical candidate is not promoted to a secure reading.\par\medskip
\par\noindent\leavevmode\phantomsection\label{GELFAND-0161-U007}\label{GELFAND-P001486:feynman-noun}\textbf{\texttt{GELFAND-0161-U007}} (feynman-noun).\nopagebreak[4]\par
The next word in the remark about Feynman resembles “эстетике”; the last joined stroke also permits “эстетику”. The following “у него” is clear.\par
Candidates: эстетике [aesthetics; dative or prepositional]; эстетику [aesthetics; accusative]. Unresolved in the diplomatic source.\par
Historical source review: Original coordinates (1125,2110,1330,2200). Frame 2408 by 3480. The complete original, line at 3x and fragment at 7x were examined. The last letter, resembling е, is visually more convincing, but the cramped ending is noted separately from the preceding verb.\par
Current source comparison: The е-like ending remains preferable, without fully excluding у. The following two words are independently legible.\par\medskip
\par\noindent\leavevmode\phantomsection\label{GELFAND-0161-U008}\label{GELFAND-P001486:names-cancellation}\textbf{\texttt{GELFAND-0161-U008}} (names-cancellation).\nopagebreak[4]\par
After “Polyakov and”, before “Zamolodchikov — Knizhnik — Polyakov”, approximately two letters are densely cancelled. A resemblance to “по” can be seen beneath the horizontal strokes.\par
Candidates: по; По. Unresolved in the diplomatic source.\par
Historical source review: Original coordinates (635,2568,755,2645). Frame 2408 by 3480. The original, line of names at 3x and cancellation at 10x were examined. The covering strokes prevent reliable recovery of the cancelled text and capitalization. The uncancelled source names and their separate underlining were preserved.\par
Current source comparison: The remaining loops do not establish either capitalization at 5x/10x. Both source order and the separately written English alphabetical order are retained.\par\medskip
\par\noindent\leavevmode\phantomsection\label{GELFAND-0161-U009}\label{GELFAND-P001486:metric-diagram}\textbf{\texttt{GELFAND-0161-U009}} (metric-diagram).\nopagebreak[4]\par
At bottom right, intersecting horizontal and vertical axes are labeled xi\_1 and xi\_2. An uneven closed contour containing g\_ab(xi) lies in the upper-right region. Thickened and pointed axis ends do not permit confident recovery of every arrow direction.\par
Candidates: horizontal axis with leftward and rightward arrows; horizontal axis with a rightward arrow; the left thickening is a pen trace; the thickened axis ends do not give confidently legible directions. Unresolved in the diplomatic source.\par
Historical source review: Original coordinates (1650,3105,2170,3420); left and vertical axis ends (1725,3160,1830,3390); right end (2010,3270,2120,3350). Frame 2408 by 3480. The complete original, drawing at 3x, left and vertical ends at 8x and right end at 10x were examined. All labels are legible. SourceFigure preserves their placement, intersecting axes and metric label inside the contour; uncertain arrowheads remain for comparison with the source.\par
Current source comparison: All three labels are clear. The partial recreation types them and retains the small endpoint hooks as strokes without assigning arrow directions. The original crop remains alongside it.\par\medskip
\subsection*{\texttt{GELFAND-P001487}}
\par\noindent\leavevmode\phantomsection\label{GELFAND-0161-U101}\label{GELFAND-P001487:top-beta-inequality}\textbf{\texttt{GELFAND-0161-U101}} (top-beta-inequality).\nopagebreak[4]\par
In “Integral (2) is def. for ...”, between beta and beta\_c is a densely corrected sign: its outer strokes form >, with additional strokes inside.\par
Candidates: beta > beta\_c; beta >= beta\_c. Unresolved in the diplomatic source.\par
Historical source review: Original 2408 by 3480, coordinates (1075,62,1270,165). Checked in the complete original, at 3x and separately at 10x. The possible earlier reading is overwritten; neither alternative is asserted.\par
Current source comparison: The outer > and extra inner strokes remain distinct at 5x/10x, but do not decide strict versus non-strict inequality.\par\medskip
\par\noindent\leavevmode\phantomsection\label{GELFAND-0161-U102}\label{GELFAND-P001487:left-margin-monogram}\textbf{\texttt{GELFAND-0161-U102}} (left-margin-monogram).\nopagebreak[4]\par
In the left margin opposite “Everything agreed!!” is a short large handwritten mark resembling joined initials.\par
Candidates: ИГ; ИМ. Unresolved in the diplomatic source.\par
Historical source review: Original 2408 by 3480, coordinates (120,1200,248,1372). Complete original, 3x and 10x; the letters and possibility of a non-letter mark were not finally distinguished.\par
Current source comparison: The joined mark still permits either monogram or a non-letter annotation. No attribution is inferred.\par\medskip
\par\noindent\leavevmode\phantomsection\label{GELFAND-0161-U103}\label{GELFAND-P001487:upper-dimension-range}\textbf{\texttt{GELFAND-0161-U103}} (upper-dimension-range).\nopagebreak[4]\par
After formula (4), the source clearly writes “This makes sense for D <= 1 and D >= 26”. The formula itself contains the square root of (D-1)(D-25).\par
Candidates: D <= 1 and D >= 26: the literal source entry. Unresolved in the diplomatic source.\par
Historical source review: Original 2408 by 3480, coordinates (330,715,1730,810); comparison formula (350,560,1750,710). Checked in the original and at 3x. Algebraic examination of the radicand also gives the boundary D=25; the manuscript number 26 was left uncorrected. The reason for choosing 26 in the source is not established.\par
Current source comparison: The enlarged views confirm both 26 in the sentence and 25 under the radical. The conjunction “and” is also retained. The real-domain calculation remains separate.\par\medskip
\par\noindent\leavevmode\phantomsection\label{GELFAND-0161-U104}\label{GELFAND-P001487:middle-cancelled-word}\textbf{\texttt{GELFAND-0161-U104}} (middle-cancelled-word).\nopagebreak[4]\par
Before “Asymptotics” is a word densely cancelled with two sloping lines; the lower parts of the letters resemble “число”.\par
Candidates: число [number]. Unresolved in the diplomatic source.\par
Historical source review: Original 2408 by 3480, coordinates (356,2195,571,2266). Checked at 3x and 10x; some letters are crossed by dense strokes, so the reading remains questioned.\par
Current source comparison: Surviving initial and final loops do not establish the whole cancelled word at 5x/10x.\par\medskip
\par\noindent\leavevmode\phantomsection\label{GELFAND-0161-U105}\label{GELFAND-P001487:right-margin-triangle-graphs}\textbf{\texttt{GELFAND-0161-U105}} (right-margin-triangle-graphs).\nopagebreak[4]\par
Two triangular drawings in the right margin. Upper: an outer triangle, an inner triangle, and three links between corresponding vertices. Thickened directional strokes occur toward the left of the inner lower edge and near the outer right vertex. Lower: an outer triangle with three segments toward an inner point; the lower-left segment ends with a small gap before that point. There are no labels.\par
Candidates: the upper thickened strokes are arrows: leftward along the inner base and rightward at the outer right corner; the thickened strokes are traces of retracing the edges rather than arrows; the lower gap is a pen break at a shared inner vertex; the lower gap intentionally distinguishes the edge endings. Unresolved in the diplomatic source.\par
Historical source review: Original 2408 by 3480, coordinates (1740,1740,2140,2300). The complete original and 3x were checked; enlarged views also received an independent repeat inspection. Clear geometry is described in SourceFigure; no direction or meaning was drawn into the doubtful strokes.\par
Current source comparison: The fresh endpoint details do not decide arrowheads versus retracing. The lower central gap is clear. The partial recreation preserves gaps and small hook remnants without inventing a central node.\par\medskip
\par\noindent\leavevmode\phantomsection\label{GELFAND-0161-U106}\label{GELFAND-P001487:bottom-cancelled-fraction}\textbf{\texttt{GELFAND-0161-U106}} (bottom-cancelled-fraction).\nopagebreak[4]\par
Before the retained right-hand part of (8), a fraction is cancelled, with a left factor resembling 2. Its numerator reads (4n+7)! or a nearby entry; n/k, plus/minus and 7/1 are crossed by strokes. Parentheses and an entry resembling n-4 are visible in the denominator; a right factorial cannot be distinguished from cancellation.\par
Candidates: 2(4n+7)!/(n-4); 2(4n+1)!/(n-4); 2(4n-7)!/(n-4); 2(4k+7)!/(k-4)!. Unresolved in the diplomatic source.\par
Historical source review: Original 2408 by 3480, coordinates (544,2840,967,3060). Checked in the complete original, at 3x and 8x; an independent repeat examination checked (535,2830,975,3070) at 8x. The transcription preserves cancellation, fraction structure and alternatives at every unclear position. The cancelled fraction was not replaced by the formula retained to its right. No complete alternative is asserted.\par
Current source comparison: The initial narrow details clipped parts of the fraction. A wider corrected 5x view and overlapping 10x details cover its numerator and denominator. Cancellation still prevents selecting a complete alternative.\par\medskip
\par\noindent\leavevmode\phantomsection\label{GELFAND-0161-U107}\label{GELFAND-P001487:bottom-final-factor}\textbf{\texttt{GELFAND-0161-U107}} (bottom-final-factor).\nopagebreak[4]\par
The last factor in the retained formula (8) has superscript k. Its base is a loop with a long hooked descender, resembling either 9 or Latin g.\par
Candidates: 9\textasciicircum{}k; g\textasciicircum{}k. Unresolved in the diplomatic source.\par
Historical source review: Original 2408 by 3480, coordinates (1390,2870,1540,3012). Original, 3x and 10x were checked; an independent repeat examination also left both readings unresolved. A literal 9\textasciicircum{}k would give an additional geometric factor relative to (7), but this mathematical consideration does not resolve the graphical ambiguity and was not used to correct the source.\par
Current source comparison: The complete loop and descender remain compatible with both candidates at 5x/10x. Asymptotic growth cannot decide the handwritten base.\par\medskip
\subsection*{\texttt{GELFAND-P001488}}
\par\noindent\leavevmode\phantomsection\label{GELFAND-0161-U110}\label{GELFAND-P001488:top-triangle-graph}\textbf{\texttt{GELFAND-0161-U110}} (top-triangle-graph).\nopagebreak[4]\par
At top left is an outer triangle and an inner triangle labeled 1 at its upper vertex, i at lower left and 2 at lower right. Corresponding outer and inner vertices are linked. Edge 1-2 is curved. On the outer left side is a thickening with a short lateral stroke resembling a downward arrowhead.\par
Candidates: the outer left side carries a downward arrow toward the lower-left vertex; the outer left thickening is overlapping ink from retracing the edge. Unresolved in the diplomatic source.\par
Historical source review: Original 2408 by 3480, coordinates (305,185,680,500). The complete original and 3x were checked; a detailed fragment was also independently re-examined. Clear labels and geometry are enumerated in SourceFigure; the doubtful arrowhead was not normalized.\par
Current source comparison: The native recreation types all three labels and curves edge 1-2. It preserves the visible small gaps near the upper inner vertex and the lower-right label, and treats the outer hook as a stroke without a direction claim.\par\medskip
\par\noindent\leavevmode\phantomsection\label{GELFAND-0161-U111}\label{GELFAND-P001488:middle-cancelled-incidence}\textbf{\texttt{GELFAND-0161-U111}} (middle-cancelled-incidence).\nopagebreak[4]\par
In the caption for G, after “Matr.” a long word resembling “инцидентности” is cancelled; “adjacency of vertices” is retained below.\par
Candidates: инцидентности [incidence]. Unresolved in the diplomatic source.\par
Historical source review: Original 2408 by 3480, coordinates (1555,1290,2010,1390). Checked in the complete original, at 3x and 6x. Some letter contours are hidden by crossing lines; the source correction and doubt about the cancelled word were preserved.\par
Current source comparison: The long cancellation still obscures enough contours to prevent accepting the complete candidate. The active adjacency wording remains clear.\par\medskip
\par\noindent\leavevmode\phantomsection\label{GELFAND-0161-U112}\label{GELFAND-P001488:middle-overdrawn-matrix-symbol}\textbf{\texttt{GELFAND-0161-U112}} (middle-overdrawn-matrix-symbol).\nopagebreak[4]\par
In the first line defining the matrix, the sign before subscript ij is repeatedly retraced: its final outer form reads G, but extra inner arcs resemble E or an earlier G.\par
Candidates: an earlier E corrected to G; a repeatedly retraced G. Unresolved in the diplomatic source.\par
Historical source review: Original 2408 by 3480, coordinates (1305,1440,1610,1560). Original, 3x and 6x were checked. G\_ij=1 was retained in the main text; a separate note accounts for uncertainty in the earlier strokes.\par
Current source comparison: The final G is clear at 5x/10x; its earlier form remains unrecovered.\par\medskip
\par\noindent\leavevmode\phantomsection\label{GELFAND-0161-U113}\label{GELFAND-P001488:middle-overdrawn-equivalence}\textbf{\texttt{GELFAND-0161-U113}} (middle-overdrawn-equivalence).\nopagebreak[4]\par
After G\_ij=0 is a double-headed double arrow crossed by several dense sloping strokes. It is unclear whether the sign has been heavily retraced or cancelled.\par
Candidates: an active equivalence sign, retraced several times; a cancelled equivalence sign. Unresolved in the diplomatic source.\par
Historical source review: Original 2408 by 3480, coordinates (1560,1560,1715,1650). The complete original, 3x and 10x were checked. The geometric arrow form is visible, but the editorial meaning of its dense strokes is not established.\par
Current source comparison: The arrow shape is visible at 5x/10x, while active versus cancelled status remains uncertain.\par\medskip
\par\noindent\leavevmode\phantomsection\label{GELFAND-0161-U114}\label{GELFAND-P001488:lower-cancelled-space-token}\textbf{\texttt{GELFAND-0161-U114}} (lower-cancelled-space-token).\nopagebreak[4]\par
Between “vector” and “space”, a short fragment is cancelled; beneath the strokes are parts of two loops or letters, possibly the beginning “пр”.\par
Candidates: пр; a short illegible letter or mathematical sign. Unresolved in the diplomatic source.\par
Historical source review: Original 2408 by 3480, coordinates (680,1968,812,2052). Original, 3x and 10x were checked. It cannot be established whether this began a repeated word or was a separate symbol; cancellation and candidates were explicitly carried over.\par
Current source comparison: Dense oblique strokes still conceal the middle. The fragment is not expanded to a word.\par\medskip
\par\noindent\leavevmode\phantomsection\label{GELFAND-0161-U115}\label{GELFAND-P001488:lower-tautological-annotation}\textbf{\texttt{GELFAND-0161-U115}} (lower-tautological-annotation).\nopagebreak[4]\par
Beneath the left O(-1) in the right-hand addition are two nearly parallel vertical curved strokes open to the right; the set \{(ell,x) | x in ell\} is at their level to the right.\par
Candidates: a vertically placed equality sign connecting O(-1) with the indicated set; a double curved annotation without an unambiguously legible mathematical sign. Unresolved in the diplomatic source.\par
Historical source review: Original 2408 by 3480, coordinates (1635,2040,2125,2195). The complete original, 3x and 5x were checked. The inclusion formula and set were confidently read; the role of the connecting annotation was described without inventing another relation.\par
Current source comparison: The two curved strokes remain distinct at 5x/10x. The visible set and inclusion are typed, but no additional equality is asserted.\par\medskip
\par\noindent\leavevmode\phantomsection\label{GELFAND-0161-U116}\label{GELFAND-P001488:bottom-cancelled-exercise-symbol}\textbf{\texttt{GELFAND-0161-U116}} (bottom-cancelled-exercise-symbol).\nopagebreak[4]\par
In the last line, between “for” and “dim V=1 and 2”, is a densely cancelled short sign with tall sloping strokes.\par
Candidates: M; V. Unresolved in the diplomatic source.\par
Historical source review: Original 2408 by 3480, coordinates (1650,3225,1760,3330). The complete original, 3x and 10x were checked. The candidates reflect surviving strokes and are not asserted; cancellation was preserved.\par
Current source comparison: The high strokes remain crossed through at 5x/10x, without establishing M or V.\par\medskip
\subsection*{\texttt{GELFAND-P001489}}
\par\noindent\leavevmode\phantomsection\label{GELFAND-0161-U201}\label{GELFAND-P001489:cancel-after-image}\textbf{\texttt{GELFAND-0161-U201}} (cancel-after-image).\nopagebreak[4]\par
To the right of “Ker v contains Im v,” is a short fragment densely cancelled with several horizontal and loop-shaped strokes. The earlier signs cannot reliably be separated from cancellation.\par
Candidates: none recorded. Unresolved in the diplomatic source.\par
Historical source review: Original 2408 by 3480; coordinates (x0,y0,x1,y1)=(1235,495,1400,566). This fragment was checked in the original, general enlargement 3x and separate enlargement 10x. The text under cancellation was not recovered; neither the number of letters nor a particular word is asserted. The body retains the cancelled fragment’s location and an Incomplete marker.\par
Current source comparison: Fresh 5x/10x views still do not yield a complete sign or word. The Incomplete marker remains distinct from an ordinary alternative reading.\par\medskip
\par\noindent\leavevmode\phantomsection\label{GELFAND-0161-U202}\label{GELFAND-P001489:cancel-before-bundle}\textbf{\texttt{GELFAND-0161-U202}} (cancel-before-bundle).\nopagebreak[4]\par
Between “Consider [source abbreviation]” and “bundle” is a cancelled short word or mathematical notation. Separate loops can be seen under the dense lines, but the entire fragment cannot be reliably read.\par
Candidates: none recorded. Unresolved in the diplomatic source.\par
Historical source review: Original 2408 by 3480; coordinates (x0,y0,x1,y1)=(527,590,686,679). The original, surrounding text at 3x and a separate 10x enlargement were examined. A possible Latin beginning is insufficient grounds to write a word; the missing content is explicitly marked Incomplete.\par
Current source comparison: The possible beginning remains insufficient after 5x/10x comparison. No whole word or letter count is supplied.\par\medskip
\par\noindent\leavevmode\phantomsection\label{GELFAND-0161-U203}\label{GELFAND-P001489:truncated-arrow-cancel}\textbf{\texttt{GELFAND-0161-U203}} (truncated-arrow-cancel).\nopagebreak[4]\par
Beneath the arrow from Lambda\textasciicircum{}k V to Lambda\textasciicircum{}(k+1) V is a small sign crossed by a sloping stroke. A separate v is clearly written above the arrow.\par
Candidates: v; k. Unresolved in the diplomatic source.\par
Historical source review: Original 2408 by 3480; checked fragment (x0,y0,x1,y1)=(665,1675,845,1820), enlargement 8x; the crossed sign itself is approximately (703,1755,759,1814). The meaning and earlier reading of the cancelled sign are not established. The v above and arrow direction were read independently; both candidates refer only to the cancelled lower sign.\par
Current source comparison: The lower sign remains unresolved at 5x/10x. The clearly separate upper v is preserved without using it to decipher the cancelled entry.\par\medskip
\par\noindent\leavevmode\phantomsection\label{GELFAND-0161-U204}\label{GELFAND-P001489:symmetry-note}\textbf{\texttt{GELFAND-0161-U204}} (symmetry-note).\nopagebreak[4]\par
In example 3, V is called “symplect. sp.”; above “2-form”, “симм.” has been added with a curved joining brace.\par
Candidates: симм. [the visible abbreviation]. Unresolved in the diplomatic source.\par
Historical source review: Original 2408 by 3480; coordinates (x0,y0,x1,y1)=(1785,2135,2110,2265), enlargement 6x, also the whole example at 3x. The strokes support “симм.”; no preceding “косо” is visible. Neighboring “symplect.” creates a possible mathematical inconsistency. No word was supplied and the form property was not corrected; the author’s intention requires historical examination.\par
Current source comparison: The added abbreviation is visible at 5x/10x, with no visible skew- prefix. Symplectic reduction is checked separately under explicitly alternating-form hypotheses.\par\medskip
\subsection*{\texttt{GELFAND-P001490}}
\par\noindent\leavevmode\phantomsection\label{GELFAND-0161-U205}\label{GELFAND-P001490:basis-abbreviation}\textbf{\texttt{GELFAND-0161-U205}} (basis-abbreviation).\nopagebreak[4]\par
In the line before e\_1,...,e\_r, the source writes “In an орто... . basis”. The abbreviation ends in dense joined strokes.\par
Candidates: ортон.; ортог.. Unresolved in the diplomatic source.\par
Historical source review: Original 2408 by 3480; coordinates (x0,y0,x1,y1)=(394,1545,605,1664), enlargement 6x; general fragment also examined at 3x. The last letter’s shape favors “ортон.”, but “ортог.” is not fully excluded. The abbreviation was not expanded and no normalization of the form was supplied.\par
Current source comparison: The last n-like shape still favors ортон. at 5x/10x. Neither a full word nor a diagonal normalization is added.\par\medskip
\par\noindent\leavevmode\phantomsection\label{GELFAND-0161-U206}\label{GELFAND-P001490:clifford-normalization}\textbf{\texttt{GELFAND-0161-U206}} (clifford-normalization).\nopagebreak[4]\par
The general relation clearly reads vv-prime + v-prime-v = inner-product(v,v-prime) h. Below it, e\_i squared = h and e\_i e\_k = - e\_k e\_i for i != k are equally clear.\par
Candidates: vv-prime + v-prime-v = inner-product(v,v-prime) h; e\_i squared = h, both as written. Unresolved in the diplomatic source.\par
Historical source review: Original 2408 by 3480; coordinates of the relations (x0,y0,x1,y1)=(425,1460,1325,1560) and (384,1662,1399,1777); checked in the original and at 3x in general fragment (255,1100,2160,1790). No factor 2 was added and the right sides were not replaced. The form normalization is not specified here; with inner-product(e\_i,e\_i)=2 the relations agree. That possible convention is not asserted as additional source text.\par
Current source comparison: Initial tight details clipped the upper relation. A corrected complete 5x strip, overlapping 10x upper details, the earlier lower details and whole original together confirm both literal relations. The normalization comparison remains conditional.\par\medskip
